\documentclass[10pt,a4paper]{amsart}
\usepackage{amsmath,amssymb,amsthm,graphicx}
% Retained preamble for the Stacks Project chapter {CH} (shared project preamble).
% The chapter file pinned-original-source.tex only carries \input{{preamble}} on its first line,
% so no preamble text lives inside the chapter itself.  The Stacks Project uses one single shared
% preamble for every chapter: the snapshot stages it three times (sources/stacks-chow/preamble.tex,
% sources/stacks-etale-cohomology/preamble.tex and focused-seeds/AGP/sources/stacks-snapshot-20261011/
% preamble.tex) and those copies are identical in size.  The block below is copied mechanically,
% line for line, from sources/stacks-chow/preamble.tex lines 14-18 and 154-217: the Xy-pic setup and
% the native theorem declarations together with the mathematical macros.  Nothing was added, removed
% or redefined; the wrapper class and the amsmath/amssymb/amsthm/graphicx base packages are supplied
% by the assembler.
% For commutative diagrams we use Xy-pic
% Chapter: stacks-spaces-more-morphisms.tex -- More on Morphisms of Spaces.
\usepackage[all]{xy}

% We use 2cell for 2-commutative diagrams.
\xyoption{2cell}

% Theorem environments.
%
\theoremstyle{plain}
\newtheorem{theorem}[subsection]{Theorem}
\newtheorem{proposition}[subsection]{Proposition}
\newtheorem{lemma}[subsection]{Lemma}

\theoremstyle{definition}
\newtheorem{definition}[subsection]{Definition}
\newtheorem{example}[subsection]{Example}
\newtheorem{exercise}[subsection]{Exercise}
\newtheorem{situation}[subsection]{Situation}

\theoremstyle{remark}
\newtheorem{remark}[subsection]{Remark}
\newtheorem{remarks}[subsection]{Remarks}

\numberwithin{equation}{subsection}

% Macros
%
\def\lim{\mathop{\mathrm{lim}}\nolimits}
\def\colim{\mathop{\mathrm{colim}}\nolimits}
\def\Spec{\mathop{\mathrm{Spec}}}
\def\Hom{\mathop{\mathrm{Hom}}\nolimits}
\def\Ext{\mathop{\mathrm{Ext}}\nolimits}
\def\SheafHom{\mathop{\mathcal{H}\!\mathit{om}}\nolimits}
\def\SheafExt{\mathop{\mathcal{E}\!\mathit{xt}}\nolimits}
\def\Sch{\mathit{Sch}}
\def\Mor{\mathop{\mathrm{Mor}}\nolimits}
\def\Ob{\mathop{\mathrm{Ob}}\nolimits}
\def\Sh{\mathop{\mathit{Sh}}\nolimits}
\def\NL{\mathop{N\!L}\nolimits}
\def\CH{\mathop{\mathrm{CH}}\nolimits}
\def\proetale{{pro\text{-}\acute{e}tale}}
\def\etale{{\acute{e}tale}}
\def\QCoh{\mathit{QCoh}}
\def\Ker{\mathop{\mathrm{Ker}}}
\def\Im{\mathop{\mathrm{Im}}}
\def\Coker{\mathop{\mathrm{Coker}}}
\def\Coim{\mathop{\mathrm{Coim}}}

% Boxtimes
%
\DeclareMathSymbol{\boxtimes}{\mathbin}{AMSa}{"02}

%
% Macros for moduli stacks/spaces
%
\def\QCohstack{\mathcal{QC}\!\mathit{oh}}
\def\Cohstack{\mathcal{C}\!\mathit{oh}}
\def\Spacesstack{\mathcal{S}\!\mathit{paces}}
\def\Quotfunctor{\mathrm{Quot}}
\def\Hilbfunctor{\mathrm{Hilb}}
\def\Curvesstack{\mathcal{C}\!\mathit{urves}}
\def\Polarizedstack{\mathcal{P}\!\mathit{olarized}}
\def\Complexesstack{\mathcal{C}\!\mathit{omplexes}}
% \Pic is the operator that assigns to X its picard group, usage \Pic(X)
% \Picardstack_{X/B} denotes the Picard stack of X over B
% \Picardfunctor_{X/B} denotes the Picard functor of X over B
\def\Pic{\mathop{\mathrm{Pic}}\nolimits}
\def\Picardstack{\mathcal{P}\!\mathit{ic}}
\def\Picardfunctor{\mathrm{Pic}}
\def\Deformationcategory{\mathcal{D}\!\mathit{ef}}
\usepackage{cleveref}
\crefname{originallemma}{original lemma}{original lemmas}
\Crefname{originallemma}{Original lemma}{Original lemmas}
\makeatletter
\newlabel{lemma-quasi-finite-separated-pass-through-finite-addendum}{{0874}{source}{}{sourcepointer-1}{}}
\newlabel{lemma-quasi-finite-separated-pass-through-finite-addendum@cref}{{[originallemma][0][]0874}{[1][0][]source}{}{}{}}
\makeatother
\crefname{originallemma}{original lemma}{original lemmas}
\Crefname{originallemma}{Original lemma}{Original lemmas}
\makeatletter
\newlabel{lemma-resolution-property-one-sheaf}{{0GV1}{source}{}{sourcepointer-2}{}}
\newlabel{lemma-resolution-property-one-sheaf@cref}{{[originallemma][0][]0GV1}{[1][0][]source}{}{}{}}
\makeatother
\crefname{originalreference}{original reference}{original references}
\Crefname{originalreference}{Original reference}{Original references}
\makeatletter
\newlabel{spaces-limits-lemma-descend-modules-finite-presentation}{{07V7}{source}{}{sourcepointer-3}{}}
\newlabel{spaces-limits-lemma-descend-modules-finite-presentation@cref}{{[originalreference][0][]07V7}{[1][0][]source}{}{}{}}
\makeatother
\crefname{originalreference}{original reference}{original references}
\Crefname{originalreference}{Original reference}{Original references}
\makeatletter
\newlabel{spaces-limits-lemma-quasi-coherent-finite-type-ideals}{{0855}{source}{}{sourcepointer-4}{}}
\newlabel{spaces-limits-lemma-quasi-coherent-finite-type-ideals@cref}{{[originalreference][0][]0855}{[1][0][]source}{}{}{}}
\makeatother
\crefname{originalreference}{original reference}{original references}
\Crefname{originalreference}{Original reference}{Original references}
\makeatletter
\newlabel{spaces-perfect-lemma-resolution-property-goes-up-affine}{{0GUU}{source}{}{sourcepointer-5}{}}
\newlabel{spaces-perfect-lemma-resolution-property-goes-up-affine@cref}{{[originalreference][0][]0GUU}{[1][0][]source}{}{}{}}
\makeatother
\crefname{originalreference}{original reference}{original references}
\Crefname{originalreference}{Original reference}{Original references}
\makeatletter
\newlabel{spaces-properties-lemma-quasi-compact-affine-cover}{{03H6}{source}{}{sourcepointer-6}{}}
\newlabel{spaces-properties-lemma-quasi-compact-affine-cover@cref}{{[originalreference][0][]03H6}{[1][0][]source}{}{}{}}
\makeatother
\title{The resolution property descends through a limit with affine transition morphisms}
\author{The Stacks Project Authors}
\date{stacks-project:spaces-more-morphisms.tex:local-snapshot-20261011; GNU FDL 1.2 or later}
\begin{document}
\maketitle
\noindent\textit{Editorial note.} Counted span is the author's own proof environment, source lines 13720-13768, immediately following the statement at 13711-13718.\par\medskip
\noindent\textit{Editorial note.} Four cross-chapter pointers (Limits of Spaces, Derived Categories of Spaces, Algebraic Spaces) are recorded in assembly.external\_references.\par\medskip
\noindent\textit{Editorial note.} The comparison of the two quotient modules is displayed inside the span; the descent lemma is then applied to it.\par\medskip
\noindent\textit{Allowed standard prerequisites (editorial).}\begin{itemize}
\item The one-sheaf criterion for the resolution property. chapter lemma lemma-resolution-property-one-sheaf, cited at proof line 13767. Turns the constructed quotient into the resolution property for the approximating space X\_i.
\item Resolution property goes up along affine morphisms. Derived Categories of Spaces, lemma spaces-perfect-lemma-resolution-property-goes-up-affine, cited at proof lines 13722-13723. Gives the easy direction: X inherits the property from X\_i.
\item Descending finitely presented modules along a limit. Limits of Spaces, lemma spaces-limits-lemma-descend-modules-finite-presentation, cited at proof lines 13761-13762. Shows the two quotient modules already agree at some stage, which is what produces the quotient on X\_i.
\item Quasi-compact affine covers and finite type ideals. Limits of Spaces, lemma spaces-limits-lemma-quasi-coherent-finite-type-ideals and Algebraic Spaces, lemma spaces-properties-lemma-quasi-compact-affine-cover, cited in the proof. Produce the affine etale presentation and the finite type ideal sheaf used to build the comparison modules.
\end{itemize}\medskip
\ifdefined\hypertarget\hypertarget{sourcepointer-1}{}\fi
\noindent\textit{Source reference.} Original lemma 0874: sources/stacks-spaces-more-morphisms/pinned-original-source.tex, source lines 7806--7806. Navigation pointer only: the counted proof cites this statement by \textbackslash{}ref as a standing input; the deduction the counted proof has to make is carried out inside the counted span, and no part of that deduction is delegated to this pointer.\par\medskip
\ifdefined\hypertarget\hypertarget{sourcepointer-2}{}\fi
\noindent\textit{Source reference.} Original lemma 0GV1: sources/stacks-spaces-more-morphisms/pinned-original-source.tex, source lines 13566--13566. Navigation pointer only: the counted proof cites this statement by \textbackslash{}ref as a standing input; the deduction the counted proof has to make is carried out inside the counted span, and no part of that deduction is delegated to this pointer.\par\medskip
\ifdefined\hypertarget\hypertarget{sourcepointer-3}{}\fi
\noindent\textit{Source reference.} Original reference 07V7: focused-seeds/geometry/S9-spaces-resolve-prop-descends/ref-tags.txt, source lines 4--4. Navigation pointer only: the counted proof cites this statement by \textbackslash{}ref as a standing input; the deduction the counted proof has to make is carried out inside the counted span, and no part of that deduction is delegated to this pointer.\par\medskip
\ifdefined\hypertarget\hypertarget{sourcepointer-4}{}\fi
\noindent\textit{Source reference.} Original reference 0855: focused-seeds/geometry/S9-spaces-resolve-prop-descends/ref-tags.txt, source lines 5--5. Navigation pointer only: the counted proof cites this statement by \textbackslash{}ref as a standing input; the deduction the counted proof has to make is carried out inside the counted span, and no part of that deduction is delegated to this pointer.\par\medskip
\ifdefined\hypertarget\hypertarget{sourcepointer-5}{}\fi
\noindent\textit{Source reference.} Original reference 0GUU: focused-seeds/geometry/S9-spaces-resolve-prop-descends/ref-tags.txt, source lines 6--6. Navigation pointer only: the counted proof cites this statement by \textbackslash{}ref as a standing input; the deduction the counted proof has to make is carried out inside the counted span, and no part of that deduction is delegated to this pointer.\par\medskip
\ifdefined\hypertarget\hypertarget{sourcepointer-6}{}\fi
\noindent\textit{Source reference.} Original reference 03H6: focused-seeds/geometry/S9-spaces-resolve-prop-descends/ref-tags.txt, source lines 7--7. Navigation pointer only: the counted proof cites this statement by \textbackslash{}ref as a standing input; the deduction the counted proof has to make is carried out inside the counted span, and no part of that deduction is delegated to this pointer.\par\medskip
\begin{lemma}
\label{lemma-resolution-property-descends}
Let $S$ be a scheme.
Let $X = \lim X_i$ be a limit of a direct system of quasi-compact
and quasi-separated algebraic spaces over $S$ with affine transition morphisms.
Then $X$ has the resolution property if and only if $X_i$ has
the resolution properties for some $i$.
\end{lemma}

\begin{proof}
If $X_i$ has the resolution property, then $X$ does by
Derived Categories of Spaces, Lemma
\ref{spaces-perfect-lemma-resolution-property-goes-up-affine}.
Assume $X$ has the resolution property.
Choose $i \in I$. We may choose an affine scheme $V_i$ and a surjective
\'etale morphism $V_i \to X_i$ 
(Properties of Spaces, Lemma
\ref{spaces-properties-lemma-quasi-compact-affine-cover}).
We may choose an embedding $j : V_i \to Y_i$
with $Y_i$ finite and finitely presented over $X_i$
(Lemma \ref{lemma-quasi-finite-separated-pass-through-finite-addendum}).
We may choose a finite type
quasi-coherent ideal $\mathcal{I}_i \subset \mathcal{O}_{Y_i}$
such that $V_i = Y_i \setminus V(\mathcal{I}_i)$
(Limits of Spaces, Lemma
\ref{spaces-limits-lemma-quasi-coherent-finite-type-ideals}).
Denote $V \to Y \to X$ the base changes of $V_i \to Y_i \to X_i$
to $X$. Denote $\mathcal{I} \subset \mathcal{O}_Y$ the pullback
of the ideal $\mathcal{I}_i$. By the easy direction of
Lemma \ref{lemma-resolution-property-one-sheaf}
there exists a finite locally free $\mathcal{O}_X$-module
$\mathcal{E}$ and a surjection $\mathcal{E} \to \pi_*\mathcal{I}$.
Note that since $\pi_i : Y_i \to X_i$ is finite and of finite presentation
we also have that $\pi : Y \to X$ is finite and of finite presentation
and that the $\mathcal{O}_{X_i}$-modules
$\pi_{i, *}\mathcal{O}_{Y_i}$ and
$\pi_{i, *}(\mathcal{O}_{Y_i}/\mathcal{I}_i)$
are of finite presentation and pullback to $X$ to give
$\pi_*\mathcal{O}_Y$ and
$\pi_*(\mathcal{O}_Y/\mathcal{I})$. Thus by Limits of Spaces, Lemma
\ref{spaces-limits-lemma-descend-modules-finite-presentation}
after increasing $i$ we can find a finite locally free
$\mathcal{O}_{X_i}$-module $\mathcal{E}_i$
and a map $\mathcal{E}_i \to \pi_{i, *}\mathcal{O}_{Y_i}$
whose base change to $X$ recovers the composition
$\mathcal{E} \to \pi_*\mathcal{I} \to \pi_*\mathcal{O}_Y$.
The pullbacks of the finitely presented $\mathcal{O}_{X_i}$-modules
$\Coker(\mathcal{E}_i \to \pi_{i, *}\mathcal{O}_{Y_i})$
and $\pi_{i, *}(\mathcal{O}_{Y_i}/\mathcal{I}_i)$
to $X$ agree as quotients of $\pi_*\mathcal{O}_Y$.
Hence by Limits of Spaces, Lemma
\ref{spaces-limits-lemma-descend-modules-finite-presentation}
we may assume that these agree, in other words that
the image of $\mathcal{E}_i \to \pi_{i, *}\mathcal{O}_{X_i}$
is equal to $\pi_{i, *}\mathcal{I}_i$.
Then we conclude that $X_i$ has the resolution property by
Lemma \ref{lemma-resolution-property-one-sheaf}.
\end{proof}

\end{document}
